\title[Finite failure schemes and quadratic pencils]{Finite failure schemes and the reconstruction of quadratic pencils}
\author{Qian Qi}
\date{September 24, 2026}
\subjclass[2020]{14M12, 14B05, 14D20, 13C40}
\keywords{Quadratic pencil, nonreduced degeneracy locus, intrinsic reconstruction, deformation, quotient stack}
\hypersetup{pdftitle={Finite failure schemes and the reconstruction of quadratic pencils, revision 149},pdfauthor={Qian Qi}}
\begin{abstract}
We study the information retained by an abstract finite failure scheme
when its ambient embedding and tensor coordinates are forgotten.
Every complex quadratic pencil, including a singular pencil, is
recovered from its ungraded local failure algebra.  Over a smooth
connected projective reduction the unmarked thickening recovers the
whole moving pencil, one constant left transformation, and the actual
source bundle without a line-twist ambiguity.  We place these inverse
statements in a relative moduli framework.  Spaces of homogeneous first
relations admit locally closed common-divisor strata with a universal
inverse to multiplication; for a determinant-power divisor the
effective first-relation stack is an explicit smooth quotient of an
open Grassmannian, without a factor-invariance hypothesis.  For a
pencil in dimension \(n\), the exact common divisor is
\(\det(T)^{n-1}\).  Its primitive residual relation intrinsically
identifies the pencil stratum.  After separately accounting for the
nonlinear substitution kernel and the ineffective right factor, its
stack is equivalent over arbitrary complex bases to
\([\Gr(2,\Sym^2 V)/\PGL(V)]\).  The equivalence retains infinitesimal
deformations and singular specializations, and its tangent complex
is the standard pencil deformation complex.  We compute the normal
directions beyond this stratum and construct projective families
through every pencil that keep the determinant radical fixed while
breaking both factor symmetries.  The first relation occurs at the
sharp uniform order \(n^2+2n-4\); the same inverse is realized on the
socle Grassmannian and on a Schubert line.  Exact coefficient
stabilizers, moving covering-map families with identical graded
vector bundles, and spectral Fitting data are also retained.
\end{abstract}
\maketitle
